\newcommand{\tipoRequerimento}{Perda por transformação}

\section{DOS FATOS E DOS DIREITOS}

\begin{enumerate}[resume]
\item A Unidade Consumidora \unidadeConsumidora{}, de titularidade do Município e devidamente enquadrada na classificação \textbf{B3 -- PODER PÚBLICO MUNICIPAL}, vem sendo faturada indevidamente com a cobrança de perdas técnicas oriundas de transformador, prática que carece inteiramente de amparo legal e regulatório para o seu grupo de faturamento.

\item A distribuidora fundamenta a referida cobrança com base no art. 305 da Resolução Normativa nº 1.000/2021 da ANEEL. Todavia, a leitura do caput do referido dispositivo deixa evidente a sua inaplicabilidade ao Grupo B:
\Citacao{
    Art. 305 A distribuidora deve adicionar aos valores medidos de energia e de demanda, ativas e reativas excedentes, a seguinte compensação de perdas para a unidade consumidora conectada do grupo A com equipamentos de medição instalados no secundário do transformador de responsabilidade do consumidor e demais usuários:
    
    I - 1\%: na conexão em tensão maior ou igual a 69 kV; ou

    II - 2,5\%: na conexão em tensão menor que 69 kV.
}

\item Ocorre que as hipóteses previstas nos incisos I e II do art. 305 são restritas e exclusivas de unidades consumidoras atendidas em \textbf{Média ou Alta Tensão (Grupo A)}. As unidades enquadradas no subgrupo B3, por definição regulatória, são atendidas em \textbf{Baixa Tensão} (tensão inferior a 2,3 kV), de modo que a medição ocorre integralmente no secundário do transformador da distribuidora.

\item Sendo o transformador de propriedade exclusiva da concessionária e ativo integrante da rede de distribuição de uso comum, as perdas técnicas decorrentes da transformação de energia não podem ser repassadas individualmente ao consumidor de baixa tensão na fatura de energia. 

\item Vale ressaltar que tais perdas técnicas de distribuição já são calculadas, reguladas e devidamente integradas pela ANEEL na composição da Tarifa de Energia (TE) e da Tarifa de Uso do Sistema de Distribuição (TUSD) aplicadas ao Grupo B, por ocasião dos processos de Reajuste e Revisão Tarifária Periódica. Portanto, a cobrança segregada na fatura configura manifesto \textit{bis in idem} (cobrança em duplicidade).

\item Diante do exposto, resta cristalino que o faturamento da referida unidade vem ocorrendo a maior, violando a estrutura tarifária vigente estabelecida pela agência reguladora e impondo ao ente público um ônus financeiro desprovido de lastro normativo.

\item Uma vez constatada a cobrança indevida, impõe-se a aplicação do art. 323 da REN nº 1.000/2021 da ANEEL, o qual assegura ao consumidor o direito à revisão do faturamento e à restituição das quantias recebidas indevidamente pela concessionária:
\Citacao{
    Art. 323. Constatado erro no faturamento, a distribuidora deve efetuar a revisão do faturamento, cobrando ou devolvendo ao consumidor a diferença apurada.

    (...)

    II - faturamento a maior: devolver ao consumidor e demais usuários, até o segundo ciclo de faturamento posterior à constatação, as quantias recebidas indevidamente nos últimos 60 ciclos de faturamento imediatamente anteriores à constatação. (Suspensos os efeitos referentes ao prazo de 60 (sessenta) ciclos, conforme DSP ANEEL 2.006, de 08.07.2024)
}

\item A manutenção desta cobrança sabidamente ilegal pela Neoenergia Pernambuco afronta diretamente o princípio da legalidade (art. 37, caput, CF/88), além de caracterizar enriquecimento ilícito por parte da concessionária, em manifesta violação ao dever de observância da modicidade tarifária.
\end{enumerate}

\section{DOS PEDIDOS}

\begin{enumerate}[resume]
\item Ante o exposto, requer-se à Concessionária:
    \begin{enumerate}
    \item que, em estrita observância às normas da Aneel e às cláusulas do contrato de concessão, promova a análise, solução e resposta conclusiva à presente reclamação no prazo máximo de 5 (cinco) dias úteis, nos termos do art. 408, I, da REN nº 1.000/2021;
    \item A imediata revisão do faturamento da Unidade Consumidora identificada, com a \textbf{exclusão definitiva} da cobrança de perdas técnicas oriundas de transformador, adequando o faturamento para observar estritamente os critérios do subgrupo B3 Poder Público;
    \item A devolução de todos os valores cobrados a maior a título de perdas por transformação nos últimos 10 (dez) anos, conforme o Despacho nº 2.006 da ANEEL, de 08 de julho de 2024, observando o prazo prescricional de 10 anos estabelecido no art. 205 do Código Civil, e respaldado pela jurisprudência administrativa e judicial vigente;
    \item Que a referida devolução dos valores seja efetuada \textbf{em dobro}, atualizados monetariamente pelo IPCA e acrescidos de juros de mora de 1\% ao mês calculados \textit{pro rata die}, conforme estabelece o § 2º do art. 323 da Resolução Normativa nº 1.000/2021 da ANEEL;
    \item Proceda com a devolução do valor integral, por meio de depósito bancário, na conta corrente do Município, conforme estabelece o § 7º, do art. 323 da REN nº 1.000/2021 da ANEEL;
    \item A expedição de resposta formal no prazo regulamentar, comunicando por escrito o resultado do deferimento desta reclamação por meio do endereço eletrônico: \emailEmpresa{}
    \end{enumerate}
\end{enumerate}